\subsection{自由李代数的导子}

命题 8. 设 X 是一个集合，M 是一个 L(X)-模，d 是 X 到 M 中的一个映射。存在且仅存在一个 L(X) 到 M 中的线性映射 D，它扩张 d 并且满足关系式：

\[
\mathrm{D} ([ a, a ^ {\prime} ]) = a. \mathrm{D} (a ^ {\prime}) - a ^ {\prime}. \mathrm{D} (a) \quad for a, a ^ {\prime} in \mathrm{L} (\mathrm{X}).\tag{20}
\]

我们借助括号定义一个以模 M × L(X) 为基础模的李代数 g

\[
[ (m, a), (m ^ {\prime}, a ^ {\prime}) ] = (a. m ^ {\prime} - a ^ {\prime}. m, [ a, a ^ {\prime} ]),\tag{21}
\]

其中 \(a, a'\) 属于 \(\mathrm{L}(\mathrm{X})\) 且 \(m, m'\) 属于 \(\mathbf{M}\) (第 I 章， § 1， no. 8)。设 \(f\) 是 \(\mathrm{L}(\mathrm{X})\) 到 \(\mathfrak{g}\) 中的同态，使得 \(f(x) = (d(x), x)\) 对一切 \(x\) （属于 \(\mathbf{X}\) ）成立；令 \(f(a) = (\mathrm{D}(a), u(a))\) （\(a\) 取遍 \(\mathrm{L}(\mathrm{X})\) ）。由公式 (21)， \(u\) 是 \(\mathrm{L}(\mathrm{X})\) 到其自身的同态；由于 \(u(x) = x\) 对每个 \(x\) （属于 \(\mathbf{X}\) ）成立，故 \(u(a) = a\) 对一切 \(a\) （属于 \(\mathrm{L}(\mathrm{X})\) ）成立，从而

\[
f (a) = (\mathrm{D} (a), a).\tag{22}
\]

于是由 (21) 与 (22)，关系式 (20) 可由 \(f([a, a']) = [f(a), f(a')]\) 得出。

反之，设 \(D'\) 是 \(\mathrm{L}(\mathrm{X})\) 到 M 中的一个映射，它满足与 (20) 类似的关系式 \((20')\) 并且扩张 d。令 \(f'(a) = (\mathrm{D}'(a), a)\) （对每个 \(a \in \mathrm{L}(\mathrm{X})\) ）；由 \((20')\) 与 \((21)\) ， \(f'\) 是 \(\mathrm{L}(\mathrm{X})\) 到 g 中的一个同态，并且在 X 上与 f 一致，从而 \(f' = f\) 且 \(D' = D\) 。

推论. \(\mathbf{X}\) 到 \(\mathrm{L}(\mathbf{X})\) 中的每个映射都可以唯一地扩张为 \(\mathrm{L}(\mathbf{X})\) 的一个导子。

当 M 取伴随表示下的 L(X) 时，关系式 (20) 表示 D 是一个导子。

